\section{Machine-Verified Proofs}\label{sec:verification}

All theoretical results in this paper are backed by a Lean~4 formalization.
The point of the artifact is not merely to restate theorems in code. It checks
the full logical chain used in the paper: local laws imply preservation;
preservation implies preference-learning equivalence; approximate local laws
imply quantitative gap bounds; theorem-backed and measurement-error extensions
recover the noisy-oracle setting; and stochastic adaptive tree objectives
inherit optimizer-transfer guarantees, including high-probability certification
versions.

\subsection{Scope}

The active development contains more than 50{,}000 lines across 123 Lean files.
It is organized into five modules:
\begin{description}[nosep]
    \item[OPT] the core Oracle-Preserving Summarization theory;
    \item[DSL] the TreePO / audit / IPW certificate layer;
    \item[CLT] supporting probability theory;
    \item[Econometrics] auxiliary identification and regression results;
    \item[ML] general supervised-learning infrastructure.
\end{description}

For this paper, the most important fact is that the formalization now reaches
all the way to the tree-based optimizer statements used in the preference
learning sections. In plain language: even when the oracle is only an
approximation to the true latent target, and even when certification is only
available on a good event, the optimizer selected from the tree objective comes
with a formally checked near-optimality guarantee.

The same is now true for the clustered-inference side of the certificate stack.
The Lean development does not stop at abstract concentration bounds. It also
formalizes the deterministic bridge from clustered confidence-interval
membership to the reported \texttt{computeDSLBound} object, together with the
calibration and oracle-error envelopes that sit on top of that interval.
Appendices~\ref{app:theorem-backed}, \ref{app:measurement-error}, and
\ref{app:clustered-inference} pull the most important of those formalized
extensions back into the paper in ordinary mathematical language.

\subsection{Axiom Policy}

The active proof stack uses one axiom: expected group loss is Lipschitz in
oracle distance. This is the mathematically precise place where the paper uses
the random-utility-model argument underlying Plackett--Luce and related
discrete-choice objectives. Everything else in the active theorem stack is
derived from explicit definitions and previously proved lemmas.

\subsection{Why This Matters}

Formal verification helps in four ways:
\begin{itemize}[nosep]
    \item \textbf{Trust}: theorems are checked by the proof assistant, not only by informal inspection.
    \item \textbf{Precision}: hidden assumptions become explicit, especially around idempotence, oracle measurement error, and adaptive-tree optimization.
    \item \textbf{Reproducibility}: anyone can rebuild the artifact from the repository.
    \item \textbf{Extensibility}: later work can reuse the existing interfaces instead of reproving the same preservation lemmas.
\end{itemize}

\subsection{Where the Details Live}

Appendix~\ref{app:lean} gives a readable summary of the formalization artifact.
The exact paper-to-Lean map, assumption crosswalk, and adaptive-tree theorem
inventory live in standalone repository documents under \texttt{lean3/docs/}.
This keeps the paper readable while still exposing the full code-level
correspondence for readers who want it.
